\begin{center}
{\sffamily\bfseries\zihao{4} 附录}
\end{center}
\addcontentsline{toc}{section}{附录}

\keypoint{附录保留核心模型实现、关键检验代码和正文未展开的补充证据。}

\newcommand{\appendixfigunit}[7]{%
  \Needspace{#1}%
  \noindent\begin{minipage}{\linewidth}
    \noindent\textbf{#2。}#3
    \par\vspace{0.25em}
    \centering
    \includegraphics[width=#4\linewidth,height=#5,keepaspectratio]{#6}
    \captionof*{figure}{#7}
  \end{minipage}
  \vspace{0.45em}
}

\subsection*{A 项目材料说明}

本文的建模、校验、绘图和论文生成均保存在同一项目目录中。附录只列出与论文结果直接相关的核心目录和关键代码，项目核心结构如下。

{\footnotesize
\begin{verbatim}
数模/
|-- data/                 原始数据、清洗数据和外部校验数据
|-- src/
|   |-- pipeline/         数据清洗与校验
|   |-- models/           传统供需基准与短期冲击模型
|   |-- calibration/      短期模型参数搜索与评价
|   |-- scenarios/        60--180 天长期三情景预测
|   |-- analysis/         敏感性、稳健性、GPR、OVX 与统计检验
|   |-- visualization/    论文正式图表生成
|-- scripts/              一键构建、R 审计与 DOCX 后处理脚本
|-- paper/                总论文 LaTeX 源码、章节和正式图表
|-- output/               模型结果、校验报告和最终 PDF/DOCX
\end{verbatim}
}

\subsection*{B 数据口径}

真实价格数据仅来自赛题附件，并经清洗后形成中文命名数据文件。短期拟合使用附件中 2026 年 3 月 3 日至 2026 年 5 月 5 日的冲突窗口；2017--2026 年完整价格序列仅用于历史同长度窗口、基准模型分布和事件极端性检验。EIA 周度 SPR、商业库存和原油产量数据、JODI 多国月度产量、库存、进出口、需求和炼厂数据、OPEC 全球供需平衡表、GPR 月度指数和 OVX 隐含波动率用于外部合理性校验，不反向替代附件真实价格。

\clearpage

\subsection*{C 短期冲击模型关键实现}

代码~\ref{lst:dynamic-recursion} 摘录短期冲击模型的关键递推逻辑，重点展示需求净扣减、商业库存自适应释放、市场感知缓冲率和价格形成层的实现方式。

\begin{paperlisting}{lst:dynamic-recursion}{python}{短期冲击模型关键递推}
def simulate_dynamic_model(event_df, assumptions, behavior):
    base_price = float(event_df.iloc[0]["pre_close"])
    previous_price = base_price
    previous_fear_excess = assumptions.fear_initial
    perceived_coverage = 0.0
    inventory_remaining = assumptions.commercial_inventory
    gap_closure_day = None
    rows = []

    for step_index, actual_row in enumerate(event_df.itertuples(index=False)):
        day_index = int((actual_row.trade_date - event_df["trade_date"].min()).days)
        elasticity = interpolate_elasticity(day_index, assumptions)

        # 需求侧：价格弹性解释一部分需求下降，题面需求收缩只扣除净额
        price_ratio = max(previous_price / base_price, 0.1)
        price_adjusted_demand = assumptions.base_demand * (price_ratio**elasticity)
        observed_decline = ramp(day_index, 0, assumptions.demand_decline_ramp_days,
                                assumptions.observed_demand_decline)
        effective_demand, demand_decline, elasticity_loss = effective_demand_after_adjustments(
            assumptions.base_demand, price_adjusted_demand, observed_decline
        )

        # 供给侧：SPR、绕道和商业库存共同缓冲供应缺口
        spr_release = ramp(day_index, assumptions.spr_delay_days,
                           assumptions.spr_ramp_days, assumptions.spr_max_release)
        route_supply = ramp(day_index, assumptions.route_start_day,
                            assumptions.route_ramp_days, assumptions.route_max_capacity)
        supply_without_inventory = (assumptions.base_supply - assumptions.supply_interruption
                                    + spr_release + route_supply)
        raw_gap = max(effective_demand - supply_without_inventory, 0.0)
        inventory_buffer, inventory_release_pressure = adaptive_inventory_buffer(
            raw_gap, previous_price, base_price, inventory_remaining, assumptions, behavior
        )
        inventory_remaining -= inventory_buffer
        residual_gap = max(effective_demand - supply_without_inventory - inventory_buffer, 0.0)

        # 市场不是立即相信所有缓冲，因此用感知覆盖率触发折价项
        gross_shortage = max(effective_demand - (assumptions.base_supply
                                                 - assumptions.supply_interruption), 0.0)
        total_buffer_supply = spr_release + route_supply + inventory_buffer
        old_coverage = perceived_coverage
        observed_coverage = np.clip(total_buffer_supply / max(gross_shortage, 1.0), 0.0, 1.5)
        perceived_coverage = update_perceived_buffer_coverage(old_coverage, observed_coverage)
        coverage_momentum = perceived_coverage - old_coverage
        if (gap_closure_day is None
                and residual_gap <= assumptions.base_demand * 0.005
                and perceived_coverage >= 0.55):
            gap_closure_day = day_index

        # 价格形成层：物理缺口、地缘风险、恐慌、不确定性与缓冲折价
        fear_excess = assumptions.fear_initial * np.exp(-assumptions.fear_decay * day_index)
        shortage_pressure = (base_price * behavior.pressure_scale
                             * (residual_gap / assumptions.base_demand)
                             / max(abs(elasticity), 0.01))
        blockade_risk_premium = (base_price * behavior.risk_weight
                                 * (assumptions.supply_interruption / assumptions.base_demand)
                                 * (1 - np.exp(-day_index / 7))
                                 * np.exp(-BLOCKADE_RISK_DECAY * day_index))
        uncertainty_premium = base_price * behavior.uncertainty_floor * (1 - np.exp(-day_index / 18))
        panic_premium = base_price * 0.45 * fear_excess
        buffer_discount = buffer_confirmation_discount(day_index, gap_closure_day,
                                                       base_price, behavior, perceived_coverage)
        relief_discount = expectation_relief_discount(day_index, base_price, behavior,
                                                      perceived_coverage, coverage_momentum)
        target_price = (base_price + shortage_pressure + blockade_risk_premium
                        + uncertainty_premium + panic_premium
                        - buffer_discount - relief_discount)
        simulated_price = previous_price + behavior.adjustment_speed * (target_price - previous_price)
        simulated_price += 2.5 * (fear_excess - previous_fear_excess)
        simulated_price = float(np.clip(simulated_price, base_price * 0.75, 180.0))

        rows.append({
            "day_index": day_index,
            "simulated_price": simulated_price,
            "supply_gap": residual_gap,
            "inventory_buffer": inventory_buffer,
            "inventory_release_pressure": inventory_release_pressure,
            "perceived_buffer_coverage_ratio": perceived_coverage,
        })
        previous_price = simulated_price
        previous_fear_excess = fear_excess

    return pd.DataFrame(rows)
\end{paperlisting}

\clearpage

\subsection*{D 参数校准模块}

代码~\ref{lst:calibration} 给出参数校准流程。本文在固定机制结构下，综合使用差分进化、局部稳定性搜索和拟合质量精修选择代表性参数。

\begin{paperlisting}{lst:calibration}{python}{参数校准与局部稳健性搜索}
def refine_with_continuous_search(event_df, base_assumptions):
    def objective(values):
        assumptions, behavior = decode_continuous_parameters(values, base_assumptions)
        simulation = dynamic.simulate_dynamic_model(event_df, assumptions, behavior)
        metrics = evaluate_simulation(simulation)

        # 综合得分同时考虑整体误差、峰值误差、末日误差和平台解释能力
        return metrics["综合得分"] + excellence_penalty(metrics)

    result = differential_evolution(
        objective,
        CONTINUOUS_PARAMETER_BOUNDS,
        seed=RANDOM_SEED,
        maxiter=LOCAL_REFINEMENT_MAXITER,
        popsize=LOCAL_REFINEMENT_POPSIZE,
        polish=True,
        workers=1,
        tol=0.002,
    )

    assumptions, behavior = decode_continuous_parameters(result.x, base_assumptions)
    simulation = dynamic.simulate_dynamic_model(event_df, assumptions, behavior)
    metrics = evaluate_simulation(simulation)
    metrics["局部精修目标值"] = float(result.fun)
    return assumptions, behavior, simulation, metrics

def calibrate(event_df, base_assumptions):
    rows, simulations = [], {}
    for candidate_id, (assumptions, behavior) in enumerate(sampled_parameter_sets(base_assumptions)):
        simulation = dynamic.simulate_dynamic_model(event_df, assumptions, behavior)
        metrics = evaluate_simulation(simulation)
        rows.append({"candidate_id": candidate_id, "candidate_source": "seeded_random",
                     **metrics})
        simulations[candidate_id] = simulation

    # 在随机搜索基础上继续做连续精修、局部稳健性搜索和拟合质量精修
    refined = refine_with_continuous_search(event_df, base_assumptions)
    stable = refine_with_local_stability_search(event_df, base_assumptions,
                                                refined[0], refined[1])
    final = refine_with_fit_quality_search(event_df, base_assumptions,
                                           stable[0], stable[1])
    return select_representative_candidates(rows, simulations, final)
\end{paperlisting}

\clearpage

\subsection*{E 中长期油价调节模型递推与政策缓冲}

代码~\ref{lst:future-recursion} 展示长期外推中 SPR 自适应释放与供需再平衡的处理。该模块避免把战略储备释放机械设为满额，而是根据剩余缺口、价格压力、时间衰减和剩余等效预算动态收缩释放强度。

\begin{paperlisting}{lst:future-recursion}{python}{中长期油价调节模型中的 SPR 自适应释放}
def adaptive_spr_release(day_index, scheduled_spr, gross_gap_before_spr,
                         previous_price, base_price, assumptions):
    if scheduled_spr <= 0:
        return 0.0, 0.0

    # 当物理缺口已基本覆盖时，SPR 不再机械满额释放
    reserve_buffer = assumptions.base_demand * GAP_CLOSURE_SHARE
    coverage_need = max(gross_gap_before_spr + reserve_buffer, 0.0)
    demand_based_release = min(scheduled_spr, coverage_need)

    stress_ratio = np.clip(gross_gap_before_spr / max(assumptions.supply_interruption, 1.0),
                           0.0, 1.0)
    price_stress = np.clip((previous_price / base_price - SPR_PRICE_STRESS_START_RATIO)
                           / SPR_PRICE_STRESS_WIDTH, 0.0, 1.0)
    minimum_policy_share = SPR_MIN_POLICY_SHARE + SPR_PRICE_POLICY_SHARE * price_stress
    stress_share = max(minimum_policy_share, stress_ratio)

    time_taper = 1.0
    if day_index > SPR_TAPER_START_DAY:
        time_taper = SPR_TAPER_FLOOR + (1 - SPR_TAPER_FLOOR) * np.exp(
            -(day_index - SPR_TAPER_START_DAY) / SPR_TAPER_DECAY_DAYS
        )

    actual_release = max(demand_based_release, scheduled_spr * stress_share * time_taper)
    actual_release = float(min(scheduled_spr, actual_release))
    return actual_release, actual_release / scheduled_spr

def simulate_future_from_prefix(prefix, future_frame, assumptions, behavior, base_price):
    spr_budget_total, spr_future_budget, spr_remaining = initialize_spr_budget(prefix, assumptions)
    inventory_remaining = float(prefix.iloc[-1]["inventory_remaining"])

    for _, row in future_frame.iterrows():
        scheduled_spr = ramp(row["day_index"], assumptions.spr_delay_days,
                             assumptions.spr_ramp_days, assumptions.spr_max_release)
        spr_release, taper_ratio = adaptive_spr_release(...)
        spr_release = min(spr_release, spr_remaining)
        spr_remaining = max(spr_remaining - spr_release, 0.0)

        # 当 SPR 接近耗尽且剩余缺口仍存在时，加入二次跳涨压力
        exhaustion_premium, exhaustion_pressure = spr_exhaustion_premium(
            spr_remaining, spr_future_budget, residual_gap, base_price, assumptions
        )
        target_price = base_price + shortage_pressure + risk_premium \
            + exhaustion_premium - buffer_discount - relief_discount
\end{paperlisting}

\clearpage

\subsection*{F 长期状态转移递推}

代码~\ref{lst:state-transition-probability} 展示长期预测中状态转移概率的递推方式。基础矩阵先给出缓和、维持、升级三类状态之间的常规切换概率，再映射到 logit 空间，由物理压力和 GPR/OVX 综合风险压力对升级倾向进行温和修正；后期缓和不再按日历时间机械触发，而由恢复证据门控。

\begin{paperlisting}{lst:state-transition-probability}{python}{长期状态转移概率递推}
def transition_probabilities(current_state, day_index, risk_constraints,
                             physical_pressure, recovery_evidence):
    base_probs = TRANSITION_MATRIX[STATE_INDEX[current_state]].copy()
    logits = np.log(np.clip(base_probs, 1e-6, 1.0))

    # 市场风险主要影响外推早中期；物理压力来自剩余缺口、SPR、库存和绕道
    early_risk_decay = np.clip(1.0 - max(day_index - 60, 0) / 150.0,
                               0.25, 1.0)
    risk_shift = risk_constraints.combined_pressure * early_risk_decay * 0.75
    physical_shift = np.clip(physical_pressure, 0.0, 1.0) * 0.80
    logits[STATE_INDEX["escalation"]] += 0.16 * risk_shift + 0.24 * physical_shift
    logits[STATE_INDEX["easing"]] -= 0.11 * risk_shift + 0.18 * physical_shift

    if current_state == "escalation":
        logits[STATE_INDEX["escalation"]] += 0.07 * risk_shift + 0.11 * physical_shift
    elif current_state == "easing":
        logits[STATE_INDEX["neutral"]] += 0.04 * risk_shift

    # 后期缓和由恢复证据门控，而不是只由日期触发
    recovery_gate = np.clip((recovery_evidence - 0.45) / 0.35, 0.0, 1.0)
    stress_brake = np.clip(1.0 - 0.75 * physical_shift, 0.0, 1.0)
    conditional_easing = 0.75 * recovery_gate * stress_brake
    if day_index >= 120:
        logits[STATE_INDEX["easing"]] += 0.26 * conditional_easing
        logits[STATE_INDEX["escalation"]] -= 0.24 * conditional_easing
    if day_index >= 150:
        logits[STATE_INDEX["easing"]] += 0.18 * conditional_easing
        logits[STATE_INDEX["neutral"]] -= 0.07 * conditional_easing
        logits[STATE_INDEX["escalation"]] -= 0.11 * conditional_easing

    exp_values = np.exp(logits - np.max(logits))
    return exp_values / exp_values.sum()
\end{paperlisting}

代码~\ref{lst:state-transition-path} 展示状态转移路径的价格递推。每条样本路径先抽取下一期市场状态，再根据状态选择对应中心情景，并叠加历史波动率约束下的扰动和状态切换跳跃，由此得到长期条件区间和尾部风险。

\begin{paperlisting}{lst:state-transition-path}{python}{状态转移路径价格递推}
def simulate_path(rng, sample_id, dates, center_price_lookup, start_price,
                  daily_sigma, transition_jump_sigma_scale, risk_constraints,
                  physical_pressure_lookup, recovery_evidence_lookup):
    state = "neutral"
    previous_price = start_price
    previous_noise = 0.0
    rows = []

    for _, base_row in dates.iterrows():
        day_index = int(base_row["day_index"])
        physical_pressure = physical_pressure_lookup.get(day_index, 0.0)
        recovery_evidence = recovery_evidence_lookup.get(day_index, 0.0)
        next_state = sample_next_state(rng, state, day_index, risk_constraints,
                                       physical_pressure_lookup,
                                       recovery_evidence_lookup)
        spec = STATES[STATE_INDEX[next_state]]
        target_price = center_price_lookup[spec.center_scenario][day_index]

        noise_scale = max(previous_price * daily_sigma
                          * spec.volatility_multiplier, 0.15)
        innovation = rng.normal(0.0, noise_scale)
        previous_noise = 0.30 * previous_noise + innovation
        transition_jump = 0.0
        if next_state != state:
            transition_jump = rng.normal(spec.transition_jump_mean,
                                         spec.transition_jump_sigma
                                         * transition_jump_sigma_scale)

        price = 0.82 * previous_price + 0.18 * target_price
        price = np.clip(price + previous_noise + transition_jump, 70.0, 145.0)
        rows.append({"sample_id": sample_id, "day_index": day_index,
                     "trade_date": base_row["trade_date"], "状态": spec.label,
                     "forecast_price": price})
        state = next_state
        previous_price = price

    return pd.DataFrame(rows)
\end{paperlisting}

\clearpage

\subsection*{G 蒙特卡洛情景预测}

代码~\ref{lst:monte-carlo} 展示长期蒙特卡洛情景树的核心思想。模型以压力指数为主轴，同时扰动供应中断、SPR 释放、绕道能力、需求弹性、风险权重和不确定性溢价，由此得到长期价格区间和尾部风险概率。

\begin{paperlisting}{lst:monte-carlo}{python}{蒙特卡洛情景树参数抽样}
def sample_parameters(rng, base_assumptions, base_behavior, historical_factors):
    raw_stress = float(rng.beta(2.1, 2.3))
    stress = float(np.clip(raw_stress + rng.normal(
        0.0, historical_factors.monte_carlo_stress_noise_sigma), 0.0, 1.0))

    supply_interruption = np.clip(1400 + 400 * stress + rng.normal(0, 35), 1400, 1800)
    spr_max_release = np.clip(700 - 500 * (stress ** 1.08) + rng.normal(0, 30), 200, 700)
    route_max_capacity = np.clip(360 - 210 * stress + rng.normal(0, 22), 150, 420)
    long_elasticity = np.clip(-0.35 + 0.25 * stress + rng.normal(0, 0.018),
                              -0.35, -0.10)

    risk_weight = np.clip(base_behavior.risk_weight * (0.74 + 0.62 * stress
                                                       + rng.normal(0, 0.06)),
                          1.65, 3.35)
    uncertainty_scale = 0.64 + 0.78 * stress + rng.normal(0, 0.06)
    uncertainty_floor = np.clip(base_behavior.uncertainty_floor * uncertainty_scale,
                                0.13, 0.36)

    assumptions = replace(base_assumptions,
                          supply_interruption=supply_interruption,
                          spr_max_release=spr_max_release,
                          route_max_capacity=route_max_capacity,
                          long_elasticity=long_elasticity)
    behavior = replace(base_behavior,
                       risk_weight=risk_weight,
                       uncertainty_floor=uncertainty_floor)
    return assumptions, behavior, {"stress_index": stress}

def run_monte_carlo(n_samples):
    forecast_frame, prefix, base_assumptions, base_behavior = load_baseline()
    historical_factors = load_historical_model_factors()
    rng = np.random.default_rng(RANDOM_SEED)
    paths, metric_rows = [], []

    for sample_id in range(1, n_samples + 1):
        assumptions, behavior, meta = sample_parameters(
            rng, base_assumptions, base_behavior, historical_factors)
        path = simulate_one_sample(forecast_frame, prefix, assumptions, behavior)
        paths.append(path)
        metric_rows.append(summarize_sample(path, prefix, meta))

    metrics = pd.DataFrame(metric_rows)
    return metrics, build_quantiles(paths), build_tail_risk(metrics)
\end{paperlisting}

\clearpage

\subsection*{H 机制消融与统计检验}

代码~\ref{lst:ablation} 用于逐项关闭模型机制，检验主结论是否依赖某一个单独设定。

\begin{paperlisting}{lst:ablation}{python}{机制消融实验设计}
VARIANTS = [
    ("完整统一模型", "题面物理机制 + 市场价格形成机制", base_variant),
    ("去除SPR释放", "关闭战略石油储备释放", no_spr),
    ("去除商业库存缓冲", "关闭库存吸收缺口能力", no_inventory),
    ("去除绕道运输", "关闭绕道运输恢复能力", no_reroute),
    ("去除需求收缩", "关闭高油价需求下降", no_demand_contraction),
    ("去除风险与不确定性溢价", "关闭地缘风险和不确定性平台", no_risk_uncertainty),
    ("去除缓冲确认折价", "关闭市场确认缓冲后的降温", no_buffer_confirmation),
    ("去除预期修复折价", "关闭中后期预期修复再定价", no_expectation_relief),
]

def run_ablation(event_df, base_assumptions, base_behavior):
    rows = []
    for name, note, variant_fn in VARIANTS:
        assumptions, behavior = variant_fn(base_assumptions, base_behavior)
        simulation = dynamic.simulate_dynamic_model(event_df, assumptions, behavior)
        metrics = evaluate_simulation(simulation)
        rows.append({
            "消融实验": name,
            "说明": note,
            "RMSE": metrics["RMSE"],
            "MAE": metrics["MAE"],
            "峰值误差": metrics["峰值误差"],
        })
    return pd.DataFrame(rows)
\end{paperlisting}

\clearpage

\subsection*{I 独立计量检验补充}

代码~\ref{lst:r-audit} 给出 R 语言独立计量检验的核心逻辑。Python 负责机制递推和情景生成，R 侧用于复核主模型相对朴素基准的误差优势、残差相关性和稳健标准误结果，从而降低单一实现口径带来的偏差风险。

\begin{paperlisting}{lst:r-audit}{r}{R 语言计量审计核心代码}
model_df <- read_csv(calibrated_path, show_col_types = FALSE) %>%
  arrange(trade_date) %>%
  mutate(
    naive_forecast = lag(actual_price),
    model_error = simulated_price - actual_price,
    naive_error = naive_forecast - actual_price
  ) %>%
  filter(!is.na(naive_forecast))

model_error <- model_df$model_error
naive_error <- model_df$naive_error

# Diebold-Mariano 检验：模型损失是否显著低于朴素上一日基准
dm_sq <- forecast::dm.test(model_error, naive_error,
                           alternative = "less", h = 1, power = 2)
dm_abs <- forecast::dm.test(model_error, naive_error,
                            alternative = "less", h = 1, power = 1)

# 校准回归与 Newey-West 稳健标准误
calibration_lm <- lm(actual_price ~ simulated_price, data = model_df)
calibration_hac <- lmtest::coeftest(
  calibration_lm,
  vcov. = sandwich::NeweyWest(calibration_lm, lag = 3,
                              prewhite = FALSE, adjust = TRUE)
)

# 残差自相关、条件异方差和波动聚集检验
lb_resid_5 <- Box.test(model_error, lag = 5, type = "Ljung-Box")
lb_sq_5 <- Box.test(model_error^2, lag = 5, type = "Ljung-Box")
arch_test <- FinTS::ArchTest(model_error, lags = 5)
\end{paperlisting}

代码~\ref{lst:sensitivity-specs} 展示敏感性分析的扰动设置。扰动对象覆盖题面物理层、政策缓冲层、需求层、行为层和长期机制透明层

\begin{paperlisting}{lst:sensitivity-specs}{python}{敏感性分析扰动参数设置}
SENSITIVITY_SPECS = [
    {"key": "supply_interruption", "参数": "供应中断量",
     "层级": "题面物理层", "扰动": [1400, 1493, 1600, 1800],
     "解释": "封锁实际造成的供应缺口。"},
    {"key": "spr_max_release", "参数": "SPR释放上限",
     "层级": "题面政策缓冲层", "扰动": [200, 450, 670, 700],
     "解释": "战略石油储备释放能力。"},
    {"key": "route_max_capacity", "参数": "绕道运输能力",
     "层级": "题面物理缓冲层", "扰动": [150, 237, 300, 420],
     "解释": "绕道和替代航线形成的额外供给能力。"},
    {"key": "long_elasticity", "参数": "长期需求弹性",
     "层级": "题面需求层", "扰动": [-0.35, -0.25, -0.18, -0.10],
     "解释": "高油价持续后需求收缩强弱。"},
    {"key": "risk_weight", "参数": "地缘风险权重",
     "层级": "市场价格形成层", "扰动": [1.85, 2.22, 2.46, 2.95],
     "解释": "市场定价地缘风险溢价的强度。"},
    {"key": "uncertainty_floor", "参数": "制度风险强度",
     "层级": "长期机制透明层", "扰动": [0.17, 0.22, 0.247, 0.30],
     "解释": "持续封锁和政策协调失败带来的不确定性。"},
]
\end{paperlisting}

\clearpage

\subsection*{J 短期冲击模型补充检验}

正文“模型建立与求解”章只保留校准、基准对比和鲁棒性摘要。附录补充给出短期模型的机制贡献、参数剖面、事件窗口、滞后平移、残差分布和局部扰动稳健性图。

\appendixfigunit{0.39\textheight}{机制贡献分解}{短期价格平台不是由单一变量决定，而是物理缺口、风险溢价、政策缓冲和预期修复共同作用。该图用于说明各机制在不同日期上的方向和相对强度。}{0.96}{0.29\textheight}{短期模型机制贡献.png}{附图 1 机制贡献分解}

\appendixfigunit{0.39\textheight}{参数剖面检验}{参数剖面图观察关键参数在最优点附近扰动时误差函数是否平滑变化。若曲线只在某一点突然变好，就可能存在单点调参风险。}{0.96}{0.30\textheight}{短期模型参数剖面图.png}{附图 2 参数一维剖面}

附图 1 和附图 2 分别回答“价格平台由哪些机制推动”和“参数是否只在单点有效”。前者给出机制方向，后者观察目标函数形态，二者共同支撑短期模型的解释性和稳定性。

\appendixfigunit{0.39\textheight}{事件窗口检验}{短期冲击窗口内部并不均匀，初始冲击、高位平台、中期回落和后期二次抬升的拟合难度不同。该图把整体误差拆到事件阶段，避免只用总 RMSE 掩盖局部表现。}{0.84}{0.23\textheight}{短期模型事件窗口检验.png}{附图 3 事件窗口检验}

\appendixfigunit{0.36\textheight}{滞后平移检验}{时间序列模型容易出现“把昨天当今天”的滞后拟合。该检验将预测曲线进行平移比较；原始位置误差最低，说明模型不是简单复制滞后价格路径。}{0.78}{0.21\textheight}{短期模型滞后平移检验.png}{附图 4 滞后平移检验}

\appendixfigunit{0.39\textheight}{残差分段分布}{残差分布用于判断误差集中在哪些阶段。中期平台阶段残差更分散，说明平台震荡和情绪再定价仍是短期模型的主要误差来源。}{0.84}{0.23\textheight}{短期残差分段分布图.png}{附图 5 残差分段分布}

\appendixfigunit{0.52\textheight}{局部扰动稳健性}{稳健性带展示当前参数邻域内的预测路径范围。若扰动后路径仍围绕真实价格平台展开，说明结论不依赖某一个极窄的参数组合。}{1.00}{0.43\textheight}{短期模型稳健性带.png}{附图 6 局部扰动稳健性}

以上六组补充检验依次对应机制来源、参数形态、阶段误差、滞后排查、残差分布和局部扰动。它们共同把正文的短期拟合结论拆成可复核证据链，说明模型优势来自机制组合，而不是单点调参或滞后复制。

\clearpage

\subsection*{K 敏感性与外生风险补充检验}

正文“模型建立与求解”章只保留参数排序、尾部风险和外生风险摘要。附录在短期模型检验之后补充长期敏感性和外生风险证据，主要用于说明扰动尺度从何而来、外部风险变量如何进入解释，以及 2026 年冲突窗口相对历史窗口是否具有特殊性。

\appendixfigunit{0.38\textheight}{历史波动参数校准}{长期蒙特卡洛扰动不应凭经验设置，因此先用 2017--2025 年同长度历史窗口估计波动、最大回撤和误差分布。该图给出扰动尺度的可追溯来源。}{0.96}{0.28\textheight}{历史波动参数校准.png}{附图 7 历史波动参数校准}

附图 7 的作用是把长期随机扰动落回历史经验区间：模型并非任意扩大波动，而是用同长度窗口的收益波动和回撤分布约束后续情景树的扰动尺度。

\appendixfigunit{0.49\textheight}{传统蒙特卡洛路径云}{该图展示不引入状态转移结构时，多参数联合扰动形成的价格路径云。它作为长期情景树的朴素随机模拟参照，用来区分普通参数扰动和状态切换风险。}{0.98}{0.32\textheight}{传统蒙特卡洛路径云图.png}{附图 8 传统蒙特卡洛路径云图}

\appendixfigunit{0.40\textheight}{GPR 地缘风险指数滞后审计}{GPR 不直接替代附件价格，也不用于同日解释同日价格。这里检验滞后风险指数与布伦特收益的关系，确认其更适合作为风险环境变量。}{0.88}{0.23\textheight}{地缘风险指数滞后审计.png}{附图 9 地缘风险指数滞后检验}

\appendixfigunit{0.36\textheight}{滞后 GPR 与月度收益}{该散点关系说明，滞后 1 月 GPR 与布伦特月度收益的线性解释力有限。因此正文只将 GPR 用作长期风险溢价和状态转移概率的约束变量。}{0.78}{0.21\textheight}{滞后GPR与布伦特收益关系.png}{附图 10 滞后 1 月 GPR 与布伦特月度收益关系}

\appendixfigunit{0.40\textheight}{OVX 隐含波动率滞后检验}{OVX 反映市场对原油期权波动率的定价预期，适合刻画不确定性和尾部风险。该图用于确认 OVX 主要约束长期波动和风险状态，而不是直接给出油价涨跌方向。}{0.96}{0.29\textheight}{OVX隐含波动率滞后检验.png}{附图 11 OVX 隐含波动率滞后检验}

\appendixfigunit{0.40\textheight}{历史同长度窗口极端性}{将 2026 年冲突窗口放回 2017--2025 年同长度窗口中比较，可以判断本题样本是否处于异常高涨幅、高波动或高预测难度区间。}{0.96}{0.31\textheight}{历史窗口极端性检验.png}{附图 12 冲突窗口历史极端性}

附图 11 和附图 12 分别从市场定价风险与历史窗口位置说明同一件事：2026 年冲突窗口不是普通平稳波动段，因此长期模型需要显式保留状态切换和尾部区间，而不能只给一条确定性均值线。

\appendixfigunit{0.42\textheight}{历史基准误差分布}{该图比较历史同长度窗口上的基准误差分布，用于说明 2026 年冲突窗口相对普通市场阶段的预测难度。}{0.96}{0.29\textheight}{R历史基准误差分布图.png}{附图 13 历史同长度窗口基准误差分布}

\appendixfigunit{0.60\textheight}{库存与供需缺口风险}{该图展示中长期外推中商业库存剩余率和剩余缺口风险的同步变化，说明悲观路径的尾部风险主要来自 SPR 等效预算耗尽、库存缓冲下降和剩余缺口同时存在。}{1.00}{0.51\textheight}{库存与供需缺口风险.png}{附图 14 库存与供需缺口风险}

上述补充图共同说明，长期外推中的随机扰动、风险变量与库存缺口并非孤立设定：历史窗口提供扰动尺度，GPR 与 OVX 提供外生风险约束，库存和剩余缺口则刻画尾部二次跳涨的物理来源。因此正文把长期结果写成条件路径、分位区间和尾部概率，而不是只给一条确定性均值预测线。

\clearpage

\subsection*{L 关键结果核查表}

\begin{table}[!htbp]
\centering
\caption*{附表 1：关键结果核查}
\begin{tabularx}{\linewidth}{lX}
\toprule
项目 & 核查结果 \\
\midrule
真实数据来源 & 真实价格来自赛题附件 CSV，外部数据只作数量级和风险约束 \\
短期窗口 & 2026 年 3 月 3 日至 2026 年 5 月 5 日，共 46 个交易日 \\
短期拟合误差 & RMSE = 3.47 美元/桶，MAE = 2.87 美元/桶，MAPE = 2.88\% \\
峰值误差 & 模拟峰值 110.55 美元/桶，实际最高收盘价 114.06 美元/桶 \\
基准对比 & 主模型优于朴素上一日基准和滚动 ARIMA 基准 \\
长期中性结果 & 第 180 天价格约 97.82 美元/桶 \\
长期悲观结果 & 第 180 天价格约 121.64 美元/桶，外推期保留二次跳涨风险 \\
尾部风险 & 状态转移情景树显示外推期最高价突破 120 美元/桶的条件概率约 85.4\% \\
最敏感因素 & 封锁风险衰减速度、SPR 释放上限、供应中断量、绕道运输能力 \\
\bottomrule
\end{tabularx}
\end{table}

